%\documentclass[journal=tosc,submission,nohyperref]{iacrtrans} % Fix for the following error I had: Package hyperxmp Error: hyperref must be loaded before hyperxmp.
\documentclass[journal=tosc,submission]{iacrtrans}
\pagestyle{plain}


%!TEX root=main.tex
% Time-stamp: <2025-06-01 15:05:37>

\usepackage[utf8]{inputenc}
\usepackage{cmap}

\usepackage[russian,english]{babel}

\usepackage{hyperref}
%\usepackage[a4paper,hmargin=1.5in,vmargin=1.5in]{geometry}
\hypersetup{colorlinks = true, urlcolor = blue, linkcolor = blue, citecolor = red}
\usepackage{cleveref}
%\usepackage{hyperxmp}

\usepackage{amsthm}
\usepackage{amsmath}
\usepackage{amssymb}
\usepackage{dsfont}
\usepackage{mathrsfs}
\usepackage[
n,
operators,
advantage,
sets,
adversary,
landau,
probability,
notions,
logic,
ff,
mm,
primitives,
events,
complexity,
asymptotics,
keys]{cryptocode}
\usepackage{awesomebox}

\usepackage{graphicx} % including pictures
\usepackage[export]{adjustbox} % use max width and max height for figures
\usepackage{xspace}

% \usepackage{mathspec}
% \setmainfont{Inria Serif}[Scale=MatchLowercase]
% \setsansfont{Inria Sans}[Scale=MatchLowercase]
% \setmathfont(Digits,Latin,Greek,Symbols)[Scale=MatchLowercase]{Inria Serif}
% \setmathrm[Scale=MatchLowercase]{Inria Serif}
%% \usepackage{mathastext}


\usepackage{algorithm}
\usepackage{algpseudocode}
\algnewcommand{\LineComment}[1]{\State{\color{gray} \(\triangleright\) #1}}
\usepackage{listings}

\usepackage{subcaption}
\usepackage{graphicx}
\usepackage{tikz}
\usetikzlibrary{decorations.pathreplacing}
\usetikzlibrary{arrows}
\usetikzlibrary{arrows.meta}
\usetikzlibrary{patterns}
\usetikzlibrary{shapes}
\usetikzlibrary{shadows}
\usetikzlibrary{calc}
\usetikzlibrary{math}
\tikzstyle{fun}=[draw,very thick,fill=white]
\tikzstyle{fork}=[shape=circle,inner sep=0pt,minimum size=3pt,fill=black]
\tikzstyle{op}=[inner sep=2pt]

\usepackage{ifthen}
\usepackage{comment}

\usepackage{booktabs}
\usepackage{multirow}

% To use monospace font in arrays easily
\usepackage{array}
\newcolumntype{t}{>{\tt}c}

\newcommand\todo[1]{{\color{red}\textbf{TODO:} #1}}

\theoremstyle{theorem}
% \newtheorem{theorem}{Theorem}
% \newtheorem{maintheorem}{Main Theorem}
% \newtheorem{proposition}{Proposition}
% \newtheorem{conjecture}{Conjecture}
% \newtheorem{observation}{Observation}
% \newtheorem{goal}{Goal}
% \newtheorem{lemma}{Lemma}
% \newtheorem{corollary}{Corollary}
% \newtheorem{openproblem}{Open Problem}

% \theoremstyle{definition}
% \newtheorem{definition}{Definition}

% \theoremstyle{remark}
% \newtheorem{remark}{Remark}

\newcommand{\qbox}[1]{{\vspace{-.3cm}\awesomebox[gray]{2pt}{\faQuestionCircle}{gray}{#1}\vspace{-.3cm}}}
\newcommand{\qboxx}[1]{\awesomebox[gray][\abShortLine]{2pt}{\faQuestionCircle}{gray}{#1}}
\newcommand*{\algcomment}[1]{{\color{gray} \text{\hspace{.5cm}$\triangleright$ #1}}}

\newcommand\adva{\ensuremath{\mathcal{A}}\xspace}
\newcommand\setX{\ensuremath{\mathcal{X}}\xspace}
\newcommand\setF{\ensuremath{\mathcal{F}}\xspace}
\newcommand\setK{\ensuremath{\mathcal{K}}\xspace}
\newcommand\ring{\ensuremath{\mathcal{R}}\xspace}
\newcommand\setY{\ensuremath{\mathcal{Y}}\xspace}
\newcommand\setI{\ensuremath{\mathcal{I}}\xspace}
\newcommand\IPM{\ensuremath{\mathcal{IP\mkern-4mu M}}\xspace}
\newcommand\class{\ensuremath{\mathcal{C}}\xspace}
\newcommand\indeg{\ensuremath{\mathsf{in}}\xspace}
\newcommand\outdeg{\ensuremath{\mathsf{out}}\xspace}
\newcommand\bit{\ensuremath{\{0,1\}}\xspace}

\newcommand\MOD{\ensuremath{\mathsf{MOD}}\xspace}
\renewcommand\AC{\ensuremath{\mathsf{AC}}\xspace}
\renewcommand\NC{\ensuremath{\mathsf{NC}}\xspace}
\newcommand\NP{\ensuremath{\mathsf{NP}}\xspace}
\newcommand\Lang{\ensuremath{\mathcal{L}}\xspace}
\newcommand\Rel{\ensuremath{\mathcal{R}}\xspace}
\newcommand\ACC{\ensuremath{\mathsf{ACC}}\xspace}
\renewcommand\TC{\ensuremath{\mathsf{TC}}\xspace}
\newcommand\PRF{\ensuremath{\mathsf{PRF}}\xspace}
\newcommand\wPRF{\ensuremath{\mathsf{wPRF}}\xspace}
\renewcommand\secpar{\ensuremath{\lambda}\xspace}
\newcommand\msk{\ensuremath{\mathsf{msk}}\xspace}
\renewcommand\key{\ensuremath{\mathsf{k}}\xspace}
\renewcommand\negl{\ensuremath{\mathsf{negl}}\xspace}
\newcommand\KeyGen{\ensuremath{\mathsf{KeyGen}}\xspace}
\newcommand\Eval{\ensuremath{\mathsf{Eval}}\xspace}

\newcommand\LL{\mathcal{L}}
\newcommand\G{\mathbb{G}}
\newcommand\Z{\mathcal{Z}}
\newcommand\N{\mathbb{N}}
\newcommand\F{\mathbb{F}}
\newcommand\ftwo{\mathbb{F}_2}
\newcommand\field[1]{\mathbb{F}_{2^{#1}}}
\newcommand\Zmod[1]{\mathbb{Z}/ {#1} \mathbb{Z}}
\newcommand\proba[1]{\mathrm{Pr}\left[ #1  \right]}
\newcommand\bigoh{\mathrm{O}}
\newcommand\trace{\mathrm{Tr}}
\newcommand\traceSub[1]{\mathrm{Tr}_{#1}}
\newcommand\charac[1]{\mathds{1}_{#1}}

\renewcommand\xor{\textsc{xor}}
\newcommand\xored{\textsc{xor}ed}
\newcommand\xors{\textsc{xor}s}
\newcommand\Tand{\textsc{And}}
\newcommand\Tands{\textsc{And}s}

\newcommand\kronecker[2]{\hat{\delta}\left( #1, #2 \right)}
\newcommand\scalarprod[2]{#1 \cdot #2}
\newcommand\walsh[1]{\mathcal{W}_{#1}}
\newcommand\ddt[1]{\delta_{#1}}
\newcommand\identity[1]{\mathcal{I}^{ #1 }}
\newcommand\hweight[1]{\mathrm{hw}\left( #1 \right)}
\newcommand\coveredby{\preccurlyeq}
\newcommand\lat[1]{\mathcal{W}_{#1}}

\renewcommand\ind[2]{\mathbf{1}_{#1}\left( #2 \right)}

\newcommand\openbutterfly[3]{\mathsf{H}^{#1}_{#2, #3}}
\newcommand\newbutterfly[1]{\mathsf{N}_{#1}}
\newcommand\linearspan[1]{\langle #1 \rangle}
\newcommand\rank{\mathrm{rank}}
\newcommand\extract[1]{\mathcal{X}_{#1}}
\newcommand\walshZeroes[1]{\mathcal{Z}_{#1}}
\newcommand\ddtZeroes[1]{\mathcal{Z}^D_{#1}}

\newcommand\spaceInput{\mathcal{V}}
\newcommand\spaceOutput{\mathcal{V}^{\perp}}
\newcommand\matSwap[1]{M_{#1}}
\newcommand\maj{\mathrm{maj}}
\newcommand\twistable[1]{#1-twistable}
\newcommand\twotwoMat[4]{\left[ \begin{array}{cc} #1 & #2 \\ #3 & #4 \end{array} \right]}

\newcommand\multInv{\mathcal{I}}

\newcommand\Lstribog{\mathsf{L}}
\newcommand\Sstribog{\mathsf{S}}
\newcommand\Pstribog{\mathsf{P}}
\newcommand\Kstribog{\mathsf{K}}
\newcommand\Lfield{\mathsf{L}^{\mathsf{F}}}

\newcommand\gen[1]{g_{#1}}
\newcommand\pmin{p_{\mathrm{min}}}
\newcommand\pkuz{p_{\mathrm{kuz}}}
\newcommand\offset{\kappa}


\newcommand\tklog{TKlog}
\newcommand\tklogmath[1]{\mathscr{T}_{#1}}
\newcommand\tkexp{TKexp}
\newcommand\tkexpmath[1]{\mathscr{T}^{-1}_{#1}}
\newcommand\fly{\log_{\alpha}^{\mathrm{FLY}}}
\newcommand\logHN{\log_{\alpha}^{\mathrm{HN}}}

\newcommand\qualDiff[1]{\mathsf{A}^{\mathrm{d}}_{ #1 }}
\newcommand\qualLin[1]{\mathsf{A}^{\ell}_{ #1 }}

\newcommand\addeq[2]{#1 \stackrel{\oplus}{\sim} #2}
\newcommand\multeq[2]{#1 \stackrel{\odot}{\sim} #2}

\newcommand\splitField[1]{\mathsf{Split}_{#1}}

% \newcommand\jacobian[2]{\mathrm{Jac}_{#2}#1}
\newcommand\jacobian[2]{\mathrm{Jac}\ \! #1 (#2)}
\newcommand\Jlin[2]{\mathrm{Jac}_{\rm lin}\ \! #1 (#2)}
\newcommand\partialDiff[2]{\frac{\partial #1}{\partial #2}}

\newcommand\permSet[1]{\mathfrak{S}_{#1}}

%% Macros Alain
%% Vectors
\newcommand{\word}[1]{\ensuremath{\boldsymbol{#1}}}
\newcommand{\av}{\word{a}}
\newcommand{\bv}{\word{b}}
\newcommand{\cv}{\word{c}}
\newcommand{\dv}{\word{d}}
\newcommand{\ev}{\word{e}}
\newcommand{\gv}{\word{g}}
\newcommand{\hv}{\word{h}}
\newcommand{\mv}{\word{m}}
\newcommand{\pv}{\word{p}}
\newcommand{\uv}{\word{u}}
\newcommand{\vv}{\word{v}}
\newcommand{\wv}{\word{w}}
\newcommand{\xv}{\word{x}}
\newcommand{\yv}{\word{y}}
\newcommand{\zv}{\word{z}}

\newcommand{\GL}[2]{\mathbf{GL}_{#1}(#2)}

\newcommand\comatrice{\mathsf{Cof}}

%% Comments
\newcommand{\yann}[1]{\textcolor{green}{\bf [Yann : #1]}}
\newcommand{\christina}[1]{\textcolor{blue}{\bf [Christina : #1]}}
\newcommand{\leo}[1]{\textcolor{red}{\bf [Léo : #1]}}
\newcommand{\geoffroy}[1]{\textcolor{magenta}{\bf [Geoffroy : #1]}}


% notations
\newcommand\aes{\textsf{AES}}
%\newcommand\distinguisher{\mathcal{D}}


\newcommand\lp[1]{\textcolor{blue!80}{\textbf{Léo:} #1}}
\newcommand\cb[1]{\textcolor{red!80}{\textbf{Christina:} #1}}
\newcommand\gc[1]{\textcolor{olive!80}{\textbf{Geoffroy:} #1}}
\newcommand\yr[1]{\textcolor{brown!80}{\textbf{Yann:} #1}}

\author{authors}
\institute{institutions}
\title{SoK: On Shallow Weak PRFs}
\subtitle{A Common Symmetric Building Block for MPC Protocols}


\begin{document}

\maketitle

\keywords{shallow weak PRFs \and alternating moduli \and Goldreich's PRG \and VDLPN \and EALPN}
\begin{abstract}
A growing number of advanced cryptographic protocols and constructions rely on symmetric primitives known as weak pseudo-random functions (wPRFs). These functions differ fundamentally from traditional PRFs: they operate in constrained models where inputs are randomly sampled and cannot be chosen by the adversary. Many such functions are instantiated in practice as shallow, non-iterated constructions.
This Systematization of Knowledge (SoK) aims to provide a unified view of shallow wPRFs by bringing together diverse constructions under a common framework. These primitives consist of a single round, and their structure differs significantly from that of classical symmetric designs. We identify and focus on four major families of such constructions: alternating moduli wPRFs, Goldreich’s pseudorandom generator family, the EA-LPN-based design, and the VDLPN construction. For each, we provide formal mathematical and algorithmic descriptions, known variants, existing cryptanalytic results, and proposed parameter sets.
Our objective is to encourage the symmetric cryptography community—particularly cryptanalysts—to analyze and estimate the security level that these primitives offer in practice. To date, their security has mostly been examined from an asymptotic point of view, and concrete cryptanalysis of fixed instances is largely missing. Given their integration into high-level cryptographic protocols, these shallow wPRFs clearly need dedicated analysis: any discovered weakness could directly impact the security of the protocols that rely on them as underlying components.
\end{abstract}

%\vfill
%\tableofcontents
%\newpage


\section{Introduction}

\yr{"While the existence or not of secure wPRFs in some complexity class could be seen as purely theoretic, it nowadays received attention for being used concretely in some advanced cryptographic protocols or other constructions. In other words, we need to fix the parameters for those constructions and the non-existence of a polynomial-time adversary gives little information about the concrete security.}

%!TODO! write introduction 

Potential name of the primitives:
\begin{itemize}
\item Shallow PRF
\item Cryptographic Shallow PRF
\item SIngle ROund Prf (SIROP)
\end{itemize}



Possible outline:
\begin{enumerate}
\item What are ``Shallow Weak PRFs''?
  \begin{enumerate}
  \item Intended Applications (MPC-friendly in a precise sense, some examples)
  \item Cost Models: minimal depth where the operations allowed vary. Severe limitation on which operation are allowed, even without considering efficiency.
  \item Security Model: weak PRF
  \end{enumerate}
\item Inventory/Zoo
  \begin{enumerate}
  \item BIPSW and variants
  \item Goldreich (including FLIP as a variant)
  \item VDLPN/EALPN
  \end{enumerate}
\item On their Cryptanalysis
  \begin{enumerate}
  \item Discussion on the security (relying on the security of the PRNG is cheating, also, how much can this be relaxed? Asymptotic vs. precise, 128 bits of security doesn't really mean 128 bits, it's more of a guesstimate)
  \item Common cryptanalysis techniques (decoding)
  \item PRF implementations and cryptanalysis challenges
  \end{enumerate}
\item Conclusion
\end{enumerate}


\subsection{Background and Notations}
\label{sec:notations}

Things we must introduce:
\begin{itemize}
\item algebraic degree
\item algebraic circuit
\item circuit depth
\end{itemize}

%Notes réunion Christina: 

%Définir MPC-friendly (toutes les opérations linéaires et les opérations non-linéaires peuvent être faites à conditions qu'ils opèrent sur un petit nombre, eg. 8 bits.) petite profondeur multiplicative. Notion principale pour toutes ces opérations.

%Grandes familles de primitives MPC: secret-sharing, tout ce qui est basé sur FHE (y compris TFHE), function secret-sharing 
%Tout cela ça vient de GGM 84




%%% Local Variables:
%%% mode: latex
%%% ispell-local-dictionary: "english"
%%% TeX-master: "main"
%%% End:

%!TEX root=main.tex

\section{Background, Notations and Definitions}
\label{sec:notations}

\subsection{Core Concepts}

Things we must introduce:
\begin{itemize}
\item algebraic degree
\item algebraic circuit
\item circuit depth
\item calculatoirement indistinguable
\item Notations $Z_n, [N]$
\end{itemize}




\subsubsection{Circuit Depth and Shallow Functions}


Weak pseudorandom functions have a broad range of applications in cryptography and beyond. Of particular relevance to this work is the notion of \emph{shallow} \wPRF. Informally, a \wPRF is said to be \emph{shallow} if it can be computed by a function of low depth in an appropriate computational model. Slightly more formally:

\paragraph{Circuits.} We are interested in \wPRF{}s that are computable by circuits within a given computational model. A circuit is a directed acyclic graph composed of two types of nodes: \emph{input nodes}, which have indegree 0, and \emph{gates}, which have indegree at least 1. Each gate with indegree~$\indeg$ and outdegree~$\outdeg$ is labeled by a function $g: \{0,1\}^\indeg \to \{0,1\}^\outdeg$.\lp{pointer que cette fois-ci on n'est pas dans $\F_p$} To evaluate a circuit~$C$ with $n$ input nodes on an input $x \in \{0,1\}^n$, we proceed as follows: each input node~$i$ is labeled with the corresponding bit~$x_i$. Then, for each gate whose $\indeg$ input wires have been assigned values $(z_1, \dots, z_\indeg)$, we compute the output $g(z_1, \dots, z_\indeg)$ and assign its $\outdeg$ bits to the outgoing wires of the gate. Gates with outdegree~0 are called \emph{output gates}, and the final output of the circuit is the bitstring composed of the values on these output wires. We define a \emph{circuit class} as a family $\mathcal{C} = \{\mathcal{C}_n\}_{n \in \mathbb{N}}$, where $\mathcal{C}_n$ denotes the set of allowed circuits with $n$ input nodes. A circuit class typically restricts the set of allowed gate types, and may impose additional structural constraints, such as bounds on indegree, outdegree, depth, size, or overall topology.
\cb{Expliquer les notions de indegree et outdegree et peut être introduire clairement la notion de depth.}
\lp{rajouter un dessin ; mettre en intro/prelimaries?}

\begin{figure}
\centering

%\usetikzlibrary{arrows.meta, positioning}

\begin{tikzpicture}[
    node distance=1cm and 1.2cm,
    input/.style={
        draw=blue, 
        thick, 
        dashed, 
        fill=blue!10, 
        minimum width=1cm, 
        minimum height=0.7cm, 
        font=\small
    },
    gate/.style={
        draw=black, 
        thick, 
        fill=gray!10, 
        minimum width=1cm, 
        minimum height=0.7cm, 
        font=\small
    },
    arrow/.style={-{Latex}, thick}
]

% Input nodes
\node[input] (x1) {\large{$x_1$}};
\node[input, above=of x1] (x0) {\large{$x_0$}};

% Gates
\node[gate, right=of x1] (g1) {\large{$g_1$}};
\node[gate] (g2) at (5,1) {\large{$g_2$}};
\node[gate, above=of g1] (g0) {\large{$g_0$}};

% Arrows to gates
\draw[arrow] (0.5,-0.125) -- (1.7,-0.125);
\draw[arrow] (0.5, 0.125) -- (1, 0.125) -- (1, 1.55) -- (1.7,1.55);
\draw[arrow] (0.5, 1.85) -- (1.7,1.85);
\draw[arrow] (2.75, 1.75) -- (3.8,1.75) -- (3.8,1.15) -- (4.5,1.15);
\draw[arrow] (2.75, 0.0) -- (3.8,0.0) -- (3.8,0.85) -- (4.5,0.85);

% Outputs from g2
\draw[arrow] (5.5,1.25) -- (6.1,1.25);
\draw[arrow] (g2) -- ++(1.1,0);
\draw[arrow] (5.5,0.75) -- (6.1,0.75);

\end{tikzpicture}

\caption{\label{fig:circuit}Graphical representation of a small circuit with two input nodes. The circuit has depth~2. The gate $g_0$ has indegree~2 and outdegree~1; $g_1$ has both indegree and outdegree equal to~1; and $g_2$ has indegree~2 and outdegree~3.}
\end{figure}


\paragraph{Shallow wPRF.} Fix a circuit class $\class$ and a function $s:\N\to\N$. We say that a pair $(\KeyGen,\Eval)$ is an $s$-shallow wPRF over $\class$ if for every $\msk$ in the range of $\KeyGen$, the function $f:x\mapsto \Eval(\msk,x)$ is computable by a depth-$s(|x|)$ circuit in $\class$. The notion of shallow wPRFs captures wPRFs that can be computed in parallel-time complexity $s$ given a sufficiently large number of computation units in the computational model defined by $\class$. \cb{On peut peut-être définir cela plus simplement ? Sinon je ne comprends pas le rôle de $s$ et ce que depth-$s(|x|)$ veut dire.}\gc{``Shallow'' veut dire ``de faible profondeur'', mais ``faible'' n'est pas formel, donc je définis ``$s$-shallow'' comme étant ``de profondeur $s$'' (depth-$s$, en anglais). Formellement, un circuit est une famille infinie de circuits (un pour chaque taille d'input possible) puisqu'il faut bien un paramètre de sécurité qui puisse grossir quelque part pour définir la sécurité asymptotique (un circuit seul est un objet de taille fixée). Les paramètres de la famille de circuit sont donc des fonctions de la taille de l'entrée. Par exemple, ``un circuit de profondeur logarithmique'' serait $\log$-shallow, puisque sur une entrée $x$ de taille $n$, il sera de profondeur au plus $\log n$.}

\bigskip
\cb{A caser cette quelque part dans cette section ou dans une section "notations" :}. 
In practice, the message length and key lengths depend on a security parameter $\lambda$. To simplify notations, we will  explicitly omit indexing all relevant parameters by $\lambda$. For instance, we will write $n$ instead of $n(\lambda)$ and $\mathcal{K}$ instead of $\mathcal{K}_{\lambda}$ for the key space.




%!TEX root=main.tex

\section{Shallow wPRFs in Higher-Level MPC Protocols}
\label{sec:applications}


The existence of pseudorandom functions computable by low-depth circuits has a long history, with strong ties to fundamental questions in complexity theory such as circuit lower bounds~\cite{STOC:RazRud94,C:MilVio12,SerVio12}, derandomization~\cite{FOCS:NisWig88,STOC:Williams13}, and learning theory~\cite{valiant1984theory,kearns1994cryptographic,C:BCGIKS21}. These connections are not the focus of this work and will not be covered, as they are typically of little relevance to cryptanalysis and symmetric cryptography.\footnote{We do note, however, that the tools and results developped in these works have significantly influenced the design of shallow wPRFs that we will cover here. Therefore, on several occasions, we will refer to some of these early works to explain the security rationale underlying some candidates.} In this work, we focus on the use of low-depth weak pseudorandom functions in practically-minded higher-level secure computation protocols (we will use the terms secure computation and MPC, short for MultiParty Computation, interchangeably). More precisely, we will be mostly concerned with cryptographic constructions that use a wPRF as an \emph{inner component} within an \emph{outer component}, where the latter is a high-end cryptographic building block (a primitive or a protocol) that uses internally the \emph{circuit description} of the wPRF (that is, we do not consider cryptographic protocols that make a ``black-box'' use of a wPRF, since such protocols are not sensitive to low-level details of the computational model where the wPRF is computed). In equivalent terms, these constructions compile a wPRF whose circuit description meets specific structural constraints into a cryptographic primitive that fulfills a set of security and efficiency requirements.


 Unlike the traditional study of shallow wPRFs, whose focus was generally on low-depth wPRFs in standard computational models of complexity theory\footnote{That is, the usual ``low-depth'' classes found in complexity-theory textbooks: $\AC^0$, $\AC^0[\oplus]$, $\ACC^0$, $\TC^0$, and related classes. They capture functions computable by constant-depth circuits operating with a certain set of gates, e.g., unbounded fan-in AND and OR gates for $\AC^0$.},  the computational models that best capture these target applications are slightly more exotic at first sight. Indeed, as will become clear next, there is a duality between the computational model in which a wPRF is implemented, and the class of functions that the outer cryptographic primitive or protocol can ``evaluate'' efficiently.

\subsection{Why Weak Pseudorandom Functions?}

Pseudorandom functions turn a short random key into a virtually unbounded number of pseudorandom strings. Advanced cryptographic protocols often consume a tremendous amount of randomness, as randomness is unavoidable to hide secret data used in complex interactions, but also to guarantee many other standard security properties (e.g., catching cheaters with random challenges, authenticating data with MACs...). As such, that pseudorandom functions can be leveraged in some of these higher-level protocols is not too surprising.

What is perhaps less natural is the focus on \emph{weak} pseudorandom functions: by definition, wPRFs require the participants to be given access to a trusted source of randomness in the first place, as pseudorandomness holds only when the inputs are truly random. However, it turns out that wherever PRFs are used as an inner component in a higher-level MPC protocol, wPRFs also suffice, for in all such applications, the inputs can be easily \emph{preprocessed} (either by the party holding them, or by all participants when they are public). When one is allowed to preprocess inputs, there is an easy way to convert any wPRF into a PRF: given an arbitrary input $x$, hash $x$ into $\Hash(x)$ using your favorite hash function (e.g., $\mathsf{SHA}$-256) before feeding it to the wPRF. This general methodology, dubbed the hash-then-evaluate paradigm in~\cite{SCN:BCEKM24}, belongs to the large family of heuristic hash-based transforms that are provably secure when the hash is modeled as a random oracle (similar to the Fiat-Shamir transform, the Fischlin transform, the Fujisaki-Okamoto transform, and many more). In addition, the work of~\cite{SCN:BCEKM24} showed that for several wPRF candidates (in fact, for essentially those covered in this SoK), instantiating $\Hash$ with a PRF whose key is publicly known yields a candidate PRF whose security guarantees are comparable to that of the original wPRF.

Therefore, if one is fine with relying on the random oracle methodology (for higher-level MPC protocols targeting real-world applications, this is usually the case) or on the security arguments outlined in~\cite{SCN:BCEKM24}, using a wPRF suffices. Furthermore, building shallow wPRFs turns out to be considerably easier than building shallow PRFs. This is a known general behavior both in theoretical cryptography\footnote{For example, PRFs cannot exist in $\AC^0$, even by using a constant number of xor and majority gates~\cite{SODA:Viola13}, while wPRFs in $\AC^0$ exist under standard assumptions such as factoring~\cite{STOC:Kharitonov93}.} and in concrete candidates design. In particular, all candidate wPRFs covered in this SoK are provably not PRFs (they are easily broken given chosen queries), and constructing shallow PRFs of a comparable complexity is an open question (or, in the case of Goldreich's PRG, is known to be impossible~\cite{SODA:Viola13,TCC:AppRay16b}).

\paragraph{On computational models and cost metrics.} There are two types of restrictions that can be enforced by the computational model, that we will losely refer to as ``hard'' and ``soft'' restrictions:
\begin{itemize}
	\item An outer scheme induces a computational model with a hard restriction when a measure of the circuits in the model (e.g., depth, gate type) cannot exceed a fixed bound. As a simple example, using a linearly-homomorphic encryption scheme (e.g., Paillier encryption~\cite{EC:Paillier99}) as outer component restricts the inner component (the primitive being homomorphically evaluated) to be computable by a circuit computing a linear combination over a ring; using instead the Boneh-Goh-Nissim encryption scheme~\cite{TCC:BonGohNis05} restricts the inner component to be computable by a \emph{depth-1} arithmetic circuit with product gates of indegree 2 (i.e., a quadratic polynomial).
	\item In other settings, the outer scheme induces a soft restriction on the circuits, where some measures (e.g., depth, gate count) can grow arbitrarily but are associated to a \emph{cost} (an efficiency metric of the final protocol, typically computation, communication, or round complexity) that increases as the measure increases. A typical example would be the use of BGV-style fully-homomorphic encryption~\cite{STOC:Gentry09,ITCS:BraGenVai12}: the inner component can be any Boolean circuit, but the computational cost (measured by noise growth, or number of necessary bootstrappings) grows with the multiplicative depth of the circuit.
\end{itemize}
Below, we cover the main target applications of these cryptographic constructions. For each target application, we provide some motivation, list the high-end cryptographic building blocks used as outer component to realize the application, and extract the computational model that captures the wPRFs that can be evaluated within each outer component, indicating whether it comes with soft or hard restrictions on the type of circuits that can be computed.

\subsection{Oblivious Pseudorandom Functions}
\label{subsec:oprf}

An oblivious pseudorandom function (OPRF)~\cite{TCC:FIPR05} for a pseudorandom function family $\PRF = (\KeyGen, \Eval)$ is a two-party protocol where one party holds a key $\msk$ in the support of $\KeyGen$, and the other party holds an input $x$ to the PRF. At the end of the protocol, the party holding $x$ should learn $\Eval(\msk,x)$ (and nothing more), while the party holding $\msk$ should not learn anything. OPRFs enjoy many applications in settings such as obliviously searching databases~\cite{TCC:FIPR05}, private set intersection and pattern matching~\cite{TCC:HazLin08,CCS:KKRT16}, password-protected secret sharing and password-authenticated key exchange~\cite{AC:JarKiaKra14}, single sign-on with privacy~\cite{EUROSP:BFHLY20}, cloud key management~\cite{CCS:JarKraRes19}, and many more---see~\cite{EUROSP:CasHesLeh22} for a recent survey on the topic.

\subsubsection{From General Secure Computation}

The simplest way to build an OPRF is to take a PRF family $\PRF = (\KeyGen, \Eval)$, and to run a secure 2-party protocol that, on input $\msk$ from one party and $x$ from the other, securely computes $\Eval(\msk,x)$. Secure $2$-party computation (2PC) refers to a general set of techniques involving two participants with private inputs $x$ and $y$ who wish to compute a public function $f(x,y)$ on their joint private inputs (with the output being revealed to one participant, or both). The two main paradigms for modern secure 2-party computation are secret-sharing-based 2PC in the preprocessing model~\cite{STOC:GolMicWig87,AC:FKOS15,CCS:KelOrsSch16,C:DPSZ12,EC:KelPasRot18} and garbled-circuit-based secure computation~\cite{FOCS:Yao86,ICALP:KolSch08,EC:ZahRosEva15,CCS:WanRanKat17b,CCS:WanRanKat17a,C:DILO22,EC:CWYY23}, the latter being typically more expensive but requiring much less rounds of communication than the former, which can make it a better option over high-latency networks. Secure computation is a very active topic: there are many available protocols, and the cost metric can vary significantly from one protocol to the other. The metrics are typically of the ``soft restriction'' type: 2PC can theoretically evaluate any function. Nevertheless, the cost is mostly driven by the following properties.
\begin{description}
\item[Free Linear Operations.] Linear operations (e.g., XORs) are typically ``for free'' in such systems, either because they can be performed locally by the participants in secret-sharing-based 2PC, or due to the free-XOR trick for garbled circuits~\cite{ICALP:KolSch08}. In contrast, the number of nonlinear gates (e.g., AND gates) matters.
\item[Low Depth.] In secret-sharing-based 2PC, the depth of the circuit (the length of the largest path from an input to an output, counted ignoring the ``free'' linear gates) is a major measure of efficiency, as it translates to a lower bound on the number of messages the participants must exchange. For example, when written as a Boolean circuit with XOR gates and AND gates, the \aes{} block cipher has depth 40~\cite{EC:ARSTZ15}. When the protocol takes place between distant machines, this incurs a significant latency: a single round trip between relatively close machines (e.g., in the same country) adds 50–100ms~\cite{SP:DKLS19}. For distant machines, latency can range from 100–260ms, up to 0.5 seconds in some cases~\cite{SP:DKLS19}. A 200ms roundtrip would result in 8s of delay for the \aes{} circuit.
\item[Correlated Randomness.] In the preprocessing model, the participants initially run a secure protocol to distribute correlated randomness among them. Different types of correlated randomness can be used to implement efficiently some ``complex'' gates. Typical examples include gates converting between arithmetic values over different fields~\cite{TCC:BIPSW18}, matrix operations~\cite{C:BCGIKS20}, or some standard non-linear functions used in machine learning such as equality tests, threshold predicates, and ReLU~\cite{TCC:BoyGilIsh19,EC:BCGGIK21}.
\end{description}
As the methods and protocols for efficiently distributing complex forms of correlated randomness evolve rapidly, it is not feasible to pinpoint a specific computational model and cost metric that would faithfully capture the evolving state-of-the-art. Nevertheless, a good rule of thumb is that secure computations shine when evaluating \emph{low-depth} circuits that minimize the number of non-linear gates, and can rely on certain types of more complex gates such as equality tests, greater-than predicates, ReLU, splines, or linear algebra, among others.

The performance gain obtained by using dedicated wPRFs to build OPRFs is significant. In Section~\ref{sec:bipsw}, we introduce the alternating moduli paradigm which is argued by its authors in~\cite[Table~2]{TCC:BIPSW18} to provide three orders of magnitude of improvement compared to OPRFs built on top of block ciphers such as \aes{}, {\sf LowMC}~\cite{EC:ARSTZ15}, or {\sf Rasta}~\cite{C:DEGLLL18}.

\subsubsection{From TFHE}

Fully-homomorphic encryption (FHE)~\cite{STOC:Gentry09} is a natural high-end outer component for building round-optimal OPRFs: using FHE, one user would encrypt $x$, let the other user compute an encryption of $\Eval(\msk,x)$ homomorphically, and decrypt the result. Unfortunately, FHE is typically considerably slower than general secure computation, and a significant effort is being invested both in optimizing FHE and in optimizing primitives for homomorphic evaluation via FHE (in the setting of ciphers, standard proposals include {\sf LowMC}~\cite{EC:ARSTZ15}, {\sf FLIP}~\cite{EC:MJSC16}, {\sf Rasta}~\cite{C:DEGLLL18}, {\sf Kreyvium}~\cite{FSE:CCFLNP16}, or {\sf Transistor}~\cite{add:Transistor}, among others). Among the best-performing FHE schemes is {\sf TFHE}~\cite{JC:CGGI20}. It features more compact ciphertexts and relatively fast linear operations, as well as a fast \emph{functional bootstrapping}~\cite{EPRINT:MicPol20,EPRINT:Joye21b} that takes only a few milliseconds.
A simplified circuit model for {\sf TFHE} is the following (also of the soft restriction type, as {\sf TFHE} can in principle evaluate any function).

\begin{description}
\item[Underlying Algebra.] The circuit operates over a small integer ring $\Z_n$ (e.g., $n=2$, $n=17$).
\item[Noise.] A quantity called ``noise'' is increased by every operation, and will lead to an incorrect decryption should it become too high. It is then crucial to prevent it from increasing too much.
\item[Programmable Bootstrap (PBS).] Bootstrapping is a ``refreshing'' operation that brings the noise back down to a given base level. In the case of {\sf TFHE}, bootstrapping is ``programmable'', meaning that it is possible to evaluate an arbitrary function $f:\Z_n\to\Z_n$ implemented by a lookup table without increasing the cost of the bootstrap. This operation is by far the most expensive, so its total number must be minimized---and it is better if those that are used can be evaluated in parallel.
\item[Free Linear Operations.] The evaluation of linear gates of arbitrary indegree is virtually ``free'' in terms of time, but it incurs a slight increase in the noise.
\end{description}

Of course, a linear function over any integer ring cannot be a wPRF, hence any wPRF in the above circuit model must contain at least one functional gate. Therefore, a natural question to answer is the following.

\qbox{Given a small integer $n$ (e.g., $n\leq 17$), are there efficient wPRF candidates computable by arithmetic circuits over $\Z_n$ with additions, multiplications by a constant, and a single layer of indegree-1 programmable bootstrapping gates? What is the smallest number of PBS gates required for a secure wPRF to exist in this model?}

Recently, the work of~\cite{EC:ADDG24} identified that a wPRF of~\cite{TCC:BIPSW18}, while not originally designed with this circuit model in mind, is especially well-suited for evaluation within {\sf TFHE}, as each wPRF evaluation can be computed using solely linear operations together with a \emph{single} layer functional bootstrapping. In fact, it is not too hard to observe that the alternative design of~\cite{TCC:BIPSW18}, denoted \emph{alternative mod-2/mod-3 wPRF} in this writeup, is a (single-bit output) wPRF that uses a \emph{single} $f$-functional gate (we cover this wPRF in Section~\ref{sec:bipsw_original}). This illustrates a general behavior, also observed in other contexts, that shallow wPRFs that optimize to the extreme in some circuit model tend to also perform well in different circuit models (in the context of block ciphers, a similar observation was made for {\sf Kreyvium}~\cite{FSE:CCFLNP16} that performs well for {\sf TFHE} in spite of being originally optimized for BGV-style FHE).


\subsection{Pseudorandom Correlation Generators and Functions}

Pseudorandom correlation generators (PCG) and pseudorandom correlation functions (PCF) are related cryptographic primitives introduced and studied in a recent line of work~\cite{CCS:BCGI18,C:BCGIKS19,CCS:BCGIKR19,C:BCGIKS20,FOCS:BCGIKS20,C:CouRinRag21,C:BCGIKR22,C:BCCD23}. They play an important role in modern secure computation protocols: secure computation considers a scenario where $n$ parties, with respective secret inputs $(x_1, \ldots, x_n)$, wish to jointly compute $f(x_1,\ldots,x_n)$ for some public function $f$ without revealing anything more than the result of the computation. Modern secure computation protocols~\cite{STOC:GolMicWig87,AC:FKOS15,CCS:KelOrsSch16,C:DPSZ12,EC:KelPasRot18} rely on \emph{correlated randomness} to achieve high concrete efficiency: random strings are sampled according to a joint distribution (``correlated'') and are securely distributed among the parties prior to the protocol. Securely generating these strings has long been a major bottleneck in these protocols. PCGs and PCFs enable multiple participants, given short correlated keys, to generate long correlated pseudorandom strings \emph{without any communication}, effectively pushing the bulk of the correlated randomness distribution to an offline computation. Below, we will describe two recent paradigms for constructing PCGs and PCFs from weak pseudorandom functions. For both paradigms, and unlike Section~\ref{subsec:oprf}, the metrics are of the ``hard restriction'' type: the class of circuits handled by the frameworks are highly specific types of very low-depth circuits that can only handle a limited set of gates.


\subsubsection{From Function Secret Sharing}
\label{subsubsec:fss}

Most constructions of PCGs and PCFs (in fact, all those cited above) follow a common template that combines a cryptographic primitive called \emph{function secret sharing} (FSS) with either a pseudorandom generator (for PCGs) or a wPRF (for PCFs). Roughly, an FSS for a function class $\setF$ is a high-level cryptographic primitive that allows sharing any function $f\in \setF$ into a pair of shares $(f_0, f_1)$ such that two parties with a common input $x$ and respective shares $f_0,f_1$ can compute $y_b = \Eval(f_b,x)$ for $b=0,1$ such that $y_0+y_1 = f(x)$. In other words, an FSS allows sharing a function $f\in\setF$ such that from the shares of $f$, two parties can obtain additive shares of $f(x)$ for any input $x$. At the heart of PCG and PCF constructions are variants of the following (semi-formal) theorem.
\begin{theorem}[Theorem~5.3 in~\cite{FOCS:BCGIKS20}]
\label{thm:pcf-ot}
Let $\ring = \{\ring_n\}_{n\in \N}$ denote a family of commutative rings. Let $\PRF=(\KeyGen,\Eval)$ denote a weak pseudorandom function and let $\setF$ denote the family of all functions $x\mapsto c\cdot\Eval(\msk,x)$ for some constant $c$ and PRF key $\msk$. If there exists an FSS for $\setF$, then there exists a PCF for the OT correlation.
\end{theorem}


The ``OT correlation'' (short for oblivious transfer correlation) is, without going into details, the most standard correlation required by secure computation protocols. Hence, the theorem above says that PCFs for useful correlations can be constructed, provided that there are wPRFs that fit into a class of functions for which there exists an FSS scheme.\footnote{More formally, we need \emph{scalar multiples} of wPRFs to fit in the class, that is, functions of the form $x\to c\cdot \Eval(\msk,x)$ for some scalar $c$.} In turn, the class of functions for which we know \emph{efficient} FSS schemes\footnote{Theoretically, there exist advanced constructions of FSS for all polynomial-time functions from strong forms of fully-homomorphic encryption~\cite{C:DHRW16} that are known to exist from the LWE assumption. However, these results are purely of theoretical interest.} is fairly limited~\cite{CCS:BoyGilIsh16}, and corresponds to the set of functions that are linear combinations of ``interval functions'', as formally defined below.

\begin{definition}
Fix a ring $\ring$ and an integer $N \in \mathbb{N}$. Let $\setI = \setI_N(\ring) \subset \ring^{N}$ denote the set of \emph{interval functions} $f_{a,b,\Delta} : [N] \to \ring$, parameterized by $a,b \in [N]$ with $a \leq b$, and $\Delta \in \ring$, and defined as:
\[
f_{a,b,\Delta}(x) = 
\begin{cases}
\Delta & \text{if } x \in [a,b], \\
0 & \text{otherwise}.
\end{cases}
\]

For a given $t \in \mathbb{N}$, define $\setF_{t,N}(\ring)$ as the set of functions of the form
\[
(x_1, \dots, x_t) \mapsto L(f_1(x_1), \dots, f_t(x_t)),
\]
where each $f_i \in \setI$ and $L : \ring^t \to \ring$ is a linear function.
\end{definition}


Equivalently, the class $\setF_{t,N}(\ring)$ can be viewed as the class of functions computable by a depth-$2$ circuit with indegree-1 comparison gates on the bottom (that output a fixed payload if the input is above, or below, a given threshold) and a $t$-ary linear combination gate on top. In the following, we will slightly abuse the definition and view $\setF_{t,N}(\ring)$ both as a class of functions and as the corresponding class of circuits. The following informal lemma summarizes what is known about the existence of efficient FSS schemes for the class $\setF$ (a more formal statement can be found in~\cite{CCS:BoyGilIsh16}).
\begin{lemma}[Informal]
\label{lemma:efficient-fss}
Assume the existence of a length-doubling pseudorandom generator. There exists an FSS for the class $\setF_{t,N}(\ring)$ where the shares $(f_0,f_1)$ of a function $f\in\setF_{t,N}(\ring)$ have size $t\cdot (\secpar\cdot \log N + \log|\ring|)$ and the cost of evaluating a share on an input $x$ is dominated by $2t(\log N + (\log|\ring|)/\secpar)$ invocations of the PRG.
\end{lemma}

In practice, the pseudorandom generator is typically instantiated using {\sf AES}, and the cost of evaluating the FSS becomes dominated by $2t(\log N + (\log|\ring|)/\secpar)$ calls to {\sf AES}, which is extremly efficient (for moderate values of $t$) using modern hardware with support for {\sf AES} instructions. 
FSS constructions for more complex classes are known, but they make a heavy use of public-key cryptography and are considerably less efficient. Therefore, past works on PCGs and PCFs have focussed on developping PRGs and wPRFs compatible with the FSS schemes promised by Lemma~\ref{lemma:efficient-fss}---that is, computable by a depth-$2$ circuit in the class $\setF$ over a suitable ring (typically, the field $\F_{2^\secpar}$, where the reader can think of $\secpar$ as being $128$):

\qbox{Are there weak pseudorandom functions computable by a circuit in $\setF_{t,N}(\F_{2^\secpar})$ for $\lambda=128$ and for some value $t \ll N$?}

In the PCF-from-FSS compiler, $N$ translates to an upperbound on the target number of samples from the correlation (for a wPRF, $N$ must be superpolynomial), and $t$ is a parameter that influences both the PCF key size and the evaluation time. The works of~\cite{FOCS:BCGIKS20,PKC:CouDuc23,C:BCGIKR22} introduced respectively the {\sf VDLPN-wPRF}, 
and the {\sf EALPN-wPRF}, as candidate solutions to the above challenge (with a value $t$ polylogarithmic in $N$). We cover these candidates in Section~\ref{sec:vdlpn} and Section~\ref{sec:ealpn} respectively.

\subsubsection{From Constrained Pseudorandom Functions}

Recently, an alternative paradigm for building PCFs has emerged in~\cite{EC:BCMPR24,AC:CDDKS24}. This paradigm replaces the outer FSS component with a constrained pseudorandom function (CPRF): a pseudorandom function equipped with a special ``constraining'' algorithm that, on the input of a constraint $C$ (a function with binary output from a target class of constraints $\class$) and a PRF key $\msk$, outputs a \emph{constrained key} $\key_C$ that allows evaluating the PRF on all inputs $x$ such that $C(x) = 1$, but leaks no information about the PRF output when $C(x)=0$. The PCFs built from this paradigm enjoy a number of appealing procedures compared to the traditional paradigm: most notably, they come with a \emph{public key setup} (a one-round protocol for securely generating the PCF keys), while efficient (one-round or multi-round) key distribution protocols for earlier PCF constructions are still an open problem. At the heart of this alternative paradigm is the following theorem:

\begin{theorem}
Let $\PRF = (\KeyGen,\Eval)$ denote a weak pseudorandom function with binary output and let $\setF$ denote the family of all functions $x\to \Eval(msk,x)$ and $x\to 1-\Eval(\msk,x)$ for some PRF key $\msk$. If there exists a CPRF for the class of constraints $\setF$, then there exists a PCF for the OT correlation.
\end{theorem}

The work of~\cite{EC:BCMPR24} additionally introduced an efficient CPRF construction where the cost of evaluating the CPRF (and therefore the PCF) is dominated by a single exponentiation over an elliptic curve (on a modern laptop and for the fastest {\sf OpenSSL} curves, this can be done in about $8\mu$s). This CPRF builds upon the seminal Naor-Reingold PRF~\cite{FOCS:NaoRei97}, and is restricted to constraints computable by an \emph{inner-product membership} (IPM) predicate.
\begin{definition}
Fix a security parameter $\secpar$, integers $B,n \in \N$, and let $g:\bit^\secpar \to [B]^n$ be a function and $S \subset \N$ be a finite set. Let $\IPM = \IPM_{B,n}(S,g)$ denote the set of \emph{inner-product membership predicates over $S$}, i.e., the functions $P_{y}:[B]^n\to \bit$ for $y \in [B]^n$ defined as follows:
\[P_{y}:x\mapsto \begin{cases} 1 \text{ if } \langle g(x),y\rangle \in S \algcomment{The inner product is computed over $\N$}\\
0 \text{ otherwise}~.
\end{cases}\]
\end{definition}
In other words, after fixing some public ``preprocessing'' function $g$ on the input, and a subset $S$, an IPM predicate is a depth-$3$ circuit with a $g$-gate at the bottom (that takes a single input $x$ and has $n$ outputs), an arbitrary $n$-ary linear combination gate in the middle, and an indegree-1 membership-test gate on top. A formal statement of the following informal lemma can be found in~\cite{EC:BCMPR24}.
\begin{lemma}[Informal]
\label{lemma:efficient-cprf}
Assume the power-DDH assumption over a group $\G$ of prime order $p$. Then there exists a CPRF for the class of constraints $\IPM = \IPM_{B,n}(S,g)$ where the size of a constrained key is $(n+|S|)\cdot \log|\G|$, and the cost of evaluating the CPRF is dominated by $n$ multiplications over $\Z_p$ and one exponentiation over $\G$.
\end{lemma}

Above, the number of multiplications can be reduced to the Hamming weight of $g(x)$. Then, to obtain efficient PCFs, it remains to find wPRFs that can be computed by an inner-product membership predicate, called IPM-wPRF in~\cite{EC:BCMPR24}:

\qbox{Are there weak pseudorandom functions computable by a circuit from $\IPM_{B,n}(S,g)$ where $B$ is some integer bound and $S$ is a small subset of integers?}

In the PCF-from-CPRF compiler, a larger $S$ translates to a larger key size (as it grows linearly with $|S|$). The work of~\cite{EC:BCMPR24} showed that two wPRF from previous works are actually IPM-wPRF: the Boneh-Ishai-Passelègue-Sahai-Wu , or alternating moduli wPRF~\cite{TCC:BIPSW18} and the Goldreich-Applebaum-Raykov wPRF~\cite{DBLP:journals/eccc/ECCC-TR00-090,TCC:AppRay16b}. We cover both candidates in Section~\ref{sec:bipsw} and Section~\ref{sec:goldreich} respectively.

\subsection{VOLE-Based Zero-Knowledge Proofs}

Zero-knowledge proofs (ZKP) are fundamental cryptographic primitives that allow proving properties about a system that depends on secret values without revealing these values. While they were mainly an object of theoretical interest after their introduction in~\cite{GolMicRac89}, four decades of research have led to a plethora of highly practical ZKP that are now routinely used in real-world systems, from authentication to anonymous blockchains and cryptocurrencies (a recent example is Google's wallet~\cite{googlewallet}).

\paragraph{Definition.} An $\NP$-language is a set of strings $\Lang\subset \bit^*$ such that there is an algorithm $\Rel_\Lang$ (the \emph{relation}) that can verify proofs of membership to $\Lang$ in polynomial-time: a string $x$ is in $\Lang$ if and only if there exists a \emph{witness} $w$, of size polynomial in $x$, such that $\Rel_\Lang(x,w)$ returns ``accept'' in time polynomial in $|x|$. A zero-knowledge proof for an $\NP$-language $\Lang$ is a two-party secure computation protocol between a prover and a verifier with the following characteristics:
\begin{itemize}
	\item Both parties have as common input a word $x \in \Lang$, and the (honest) prover additionally holds a \emph{witness} for $x$---that is, a string $w$ such that $\Rel_\Lang(x,w)=1$.
	\item At the end of the interaction, the verifier outputs 1 (``accept'') or 0 (``reject'').
	\item An honest prover (with a witness $w$) should always cause the verifier to accept. However, if $x\notin\Lang$ (hence no valid witness $w$ exists), no prover can possibly cause the verifier to output 1 with non-negligible probability.
	\item Eventually, the proof is \emph{zero-knowledge}, meaning that it leaks no information about $w$ (beyond the fact that $x\in \Lang$). This is formalized by requiring that the view of the verifier (the transcript of its interaction with the prover) could have been simulated without the witness $w$, and the simulation is indistinguishable from the real transcript (hence, in essence, the verifier cannot tell the transcript apart from a transcript computed independently of $w$).
\end{itemize}

\paragraph{Efficient zero-knowledge proof systems.} Modern zero-knowledge proof systems require tradeoffs and compromises. Typically, they will handle restricted types of algebraic statements~\cite{C:Schnorr89,EC:GroSah08}, or be fully generic and optimal on almost any aspect, but incur an often prohibitive computational overhead on the prover side~\cite{EC:BCCGP16,EC:Groth16,CCS:AHIV17,SP:BBBPWM18,ICALP:BBHR18,EC:BCRSVW19,CCS:MBKM19,C:Setty20,EC:CBBZ23} (the citations are a short sample from the vast SNARK/STARK ecosystem of proofs). Building on the seminal work of~\cite{STOC:IKOS07}, a recent line of work has put forth \emph{VOLE\footnote{Short for Vector Oblivious Linear Evaluation: a mechanism that distributes vectors $\vec u, \vec v$ to the prover and $\Delta, \vec u+\Delta\cdot \vec v$ to the verifier, where $\Delta$ acts as a MAC to ``authenticate'' the vector $\vec v$.}-based zero-knowledge} as a promising middle ground~\cite{CCS:BCGI18,SP:WYKW21,ITC:DitIshOst21,CCS:YSWW21,C:BMRS21}: VOLE-based ZKP can handle arbitrary statements, have reasonable proof sizes (though higher than SNARKs/STARKs for large statements), and small prover costs, confering them an overall highly competitive scalability. VOLE-based ZKP are deployed in real-world applications (e.g.\ TLSNotary~\cite{tlsnotary}), and in commercial solutions (such as {\sf zkPass}~\cite{zkpass} or {\sf zkTLS}~\cite{primus}).

A typical ``pain point'' of zero-knowledge proofs are \emph{mixed statements}: proofs of statements that mix arithmetics over different structures (a typical example is neural network training, that often mixes arithmetic over large fields with Boolean operations). Mixed statements have been the focus of a significant attention, both in general secure computation~\cite{NDSS:DemSchZoh15,SP:GYKW24,EC:BCGGIK21} and in zero-knowledge proofs~\cite{C:ChaGanMoh16,CCS:BBMRS21,EPRINT:OKMZ24,DBLP:conf/eurocrypt/AgarwalBBS25}. A recent work~\cite{DBLP:conf/eurocrypt/AgarwalBBS25} observed that shallow wPRFs provide an ideal tool to handle mixed arithmetic in VOLE-based zero-knowledge proofs. At the high level, the key idea is the following: in VOLE-based zero-knowledge, the secret witness is authenticated to the verifier via an information-theoretic homomorphic MAC. Previous works had observed that when the secret witness must be used over two different structures, e.g., the field $\F_2$ and a larger field $\F_p$, it suffices to have pairs of identical bits MAC-authenticated both in $\F_2$ and $\F_p$ to securely ``switch'' between the two structures. The work of~\cite{DBLP:conf/eurocrypt/AgarwalBBS25} observes that an efficient way to achieve this is by first authenticating a short wPRF key over both fields and proving the equality. This is a mixed statement, so it can be expensive, but it is crucially independent of the size of the full statement since the key is short: it will get amortized away. Then, identical \emph{pseudorandom} bits are generated over both fields by evaluating the wPRF and proving (separately over each field) the correct evaluation. To make this procedure efficient, the key is to find a wPRF that admits short proofs of correct evaluation \emph{both over $\F_2$ and $\F_p$}.

\paragraph{Cost model.} The basic cost model of modern VOLE-based zero-knowledge proofs is the usual ``pay-per-product'' MPC cost model, inherited from {\sf Quicksilver}~\cite{CCS:YSWW21} (the underlying proof system used in these works): when proving that an arithmetic circuit $C$ (with additions and multiplications over a field $\F$) evaluates to a certain value on a secret witness $w$, all linear operations (additions, multiplications by a constant) are for free (they do not incur any overhead in the proof size), while multiplications increase the proof size. Therefore, the wPRF must be of the standard ``MPC-friendly'' type: its circuit description must minimize the number of multiplications.

In the setting of mixed statements, however, a subtlety arises: given a statement mixing operations over two fields $(\F,\F')$, the same wPRF must be computed by an arithmetic circuit over each of these fields. One must therefore be careful, as a low multiplication-gate count over one field does not typically translate to a low multiplication-gate count over the other field---in fact, the opposite is often true: linear operations over a field $\F_p$ are highly nonlinear over a field $\F_q$ when $p,q$ are coprime~\cite{Raz87,Smo87} (this behavior is actually a crucial component of the alternating moduli paradigm, that we cover in Section~\ref{sec:bipsw}). However, when $\F,\F'$ are prime order fields, a simple observation is that it suffices to have a wPRF computable by a circuit with a small number of multiplications \emph{over the integers}\footnote{The circuit operates directly on the input bits of the wPRF, and must output integers that are guaranteed to be bits as well.}: the natural arithmetization of this circuit (replacing additions and multiplications over the integers with the field operations) yields a valid arithmetic circuit for the wPRF over any prime-order field.
 In turn, any wPRF computable by a $d$-\emph{local} Boolean circuit (i.e., such that each output bit depends on at most $d$ input bits) can be arithmetized into an arithmetic circuit over the integers with at most $d$ multiplications. Combining these two observations, we get the following target cost model.
 \begin{description}
\item[Underlying Algebra.] The circuit operates over the Boolean field $\F_2$.
\item[Locality.] The circuit is ``$d$-local'' for some small quantity $d$: every output bit depends on at most $d$ input bits.
\end{description}

The question below summarizes the quest for the best possible candidate:

\qbox{Given a target bound $n(\secpar)$ on the key size (for a security parameter $\secpar$) and a target number $N(\secpar)$ of samples, what is the smallest $d$ such that there exist weak pseudorandom functions with key size $n(\secpar)$ computable by a $d$-local Boolean circuit?}

The study of local PRGs and wPRFs has a rich history in cryptography. We cover the main candidates in Section~\ref{sec:goldreich}.




\input{bipsw}
\section{Goldreich's PRG}
\label{sec:goldreich}

Twenty-five years ago, Goldreich posed the question of whether it is possible to securely construct a very simple function that is hard to invert~\cite{DBLP:journals/eccc/ECCC-TR00-090}. By \emph{very simple}, we mean a function whose description is transparent and can be expressed using a concise, closed-form formula. This question led to the design of a minimalistic and elegant one-way function. Despite its simplicity, this function is conjectured to offer strong cryptographic guarantees, and has, over the years, become a central object of study in the theory of cryptography and complexity.  %While the function's structure is very simple, its security also relies on access to a large collection of public, uniformly random values. \cb{Il faut dire deux mots de plus sur ce que cette dernière phrase implique.}\yr{C'est bizarre de l'expliquer ici et surtout avant de décrire la fonction ? Je parle juste de ce qu'on s'est dit lundi: weak nonce et uniformly subsets, qui si on les a pas, la construction ne marche plus.}


\subsection{Original Reference}

Goldreich's proposal for constructing a one-way function is based on a very simple and elegant idea. Let $ n, m \in \mathbb{N}$, and let $ (\sigma_i)_{1 \leq i \leq m}$   be a sequence of ordered subsets of $\{1, \dots, n\}$, each of size $ d(n)$. The parameter $d(n)$ is referred to as the \emph{locality} of the construction. Let $P : \{0,1\}^{d(n)} \to \{0,1\}$ be a fixed Boolean predicate  (i.e., a function on $d(n)$ input bits).

For any input $x \in \{0,1\}^n$, and for each $i \in \{1, \dots, m\}$, define  $x[\sigma_i] \in \{0,1\}^{d(n)}$   as the projection of $x$   onto the coordinates specified by $ \sigma_i$  :
\[
x[\sigma_i] = \left( x_{\sigma_i(1)},\; x_{\sigma_i(2)},\; \dots,\; x_{\sigma_i(d(n))} \right).
\]

The candidate one-way function $F: \{0,1\}^n \to \{0,1\}^m$   is then defined as:
\[
F(x) = \left( P(x[\sigma_1]),\; P(x[\sigma_2]),\; \dots,\; P(x[\sigma_m]) \right).
\]

We define the \emph{stretch} $\mathsf{s} \geq 1$ as the integer in the expression  $m = \mathcal{O}(n^{\mathsf{s}})$. A fundamental question in the study of Goldreich’s construction is: how large can the stretch be for a fixed locality parameter $d$, while still maintaining security?

This function is specified by a fixed predicate and a public collection of access patterns. It can be instantiated to yield various cryptographic primitives, such as one-way functions (when $m = n$), \cb{on l'a déjà défini comme une one-way function plus en haut, je ne comprends donc pas la phrase et pourquoi $m = n$?} pseudo-random functions~\cite{TCC:AppRay16b} or pseudorandom generators with polynomial stretch \cb{je ne comprends pas ce que polynomial stretch veut dire},  as the existence of local pseudorandom generators follows from the existence of one-way random local functions~\cite{STOC:Applebaum12}.



Regarding security, Goldreich's construction is usually used as a pseudorandom generator (PRG). That is, for a uniformly random seed $x \in \{0,1\}^n$, the output sequence $y = F(x) \in \{0,1\}^m$ should be computationally indistinguishable from a uniformly random string of the same length.

While this construction may initially seem purely theoretical, its structure and constraints provide valuable insights for practical cryptographic applications. In particular, the \emph{locality} parameter $d(n)$ serves as a measure of the circuit's arithmetic depth or of the complexity of the circuit implementing the functions, and plays an important role in balancing efficiency and security.

%The stretch $\mathsf{s} \geq 1$ is defined as the polynomial extension of the generator: $m = \mathcal{O}(n^{\mathsf{s}})$. One of the fundamental questions is how big can the stretch can be for a given locality $d$ by remaining secure ?

%About security, Goldreich's is a pseudorandom generator, in other words it is secure as soon as the output sequence $y=F(x)$ is indistinguishable from a random one without naturally the knowledge of the seed $x$. \yr{Est-ce que ça suffit ou bien on veut quelque chose de bien plus formel ?}

%While this construction may at first appear purely theoretical, its security properties offer valuable insights for practical applications. In particular, the \emph{locality} parameter $d(n)$ serves as a useful measure of the maximum arithmetic depth or complexity of the circuit implementing the function.

One of the most compact instantiations of this proposal uses locality $d = 5$ and the predicate
\[
P_5(x_1, x_2, x_3, x_4, x_5) = x_1 + x_2 + x_3 + x_4 x_5.
\]



More generally, natural candidates for the predicate $P$  in Goldreich’s construction include the classes of \textsf{XOR-MAJ} and \textsf{XOR-THR} functions. These correspond to predicates defined as the XOR of the output of a linear function and either a majority or a threshold function applied to independent input bits. 

These predicates are the two most popular choices, as they offer both high algebraic immunity~\cite{C:Courtois03,EC:CouMei03} and a high resiliency order~\cite{DBLP:journals/tit/Siegenthaler84}. For the \textsf{XOR-MAJ} and \textsf{XOR-THR} functions in particular, these cryptographic criteria have been thoroughly analyzed in~\cite{DBLP:journals/ccds/Meaux21,DBLP:journals/dam/Meaux22}.

Another interesting, though theoretical, question is whether there exists a predicate that simultaneously achieves optimal resiliency order and algebraic immunity for any given locality $d$~\cite{DBLP:journals/dcc/DupinMR23}.
\yr{J'ai mis ça là, mais avant c'était après les cryptanalyses qui expliquent, ça peut se descendre.}
 

\medskip

More formally, let $\mathsf{w}(x)$  denote the Hamming weight of $x \in \{0,1\}^t$ , where $t$  is the input length.

\paragraph{Threshold function.} For a given threshold $\mathsf{th} \in [1,t+1] $ , the \emph{threshold function} $\mathsf{THR}_{\mathsf{th}} : \{0,1\}^t \to \{0,1\}$  is defined as:
\[
\mathsf{THR}_{\mathsf{th}}(x) = 
\begin{cases}
0 & \text{if } \mathsf{w}(x) < \mathsf{th}, \\
1 & \text{otherwise}.
\end{cases}
\]

\paragraph{Majority function.} The \emph{majority function} $\mathsf{MAJ} : \{0,1\}^t \to \{0,1\}$  is the special case where the threshold is set to half the input length (rounded up):
\[
\mathsf{MAJ}(x) = \mathsf{THR}_{\lfloor t/2 \rfloor}(x)\,.
\]


%More generally the \textit{a priori} good candidates for predicates are the $\mathsf{XOR-MAJ}$ and $\mathsf{XOR-THR}$ predicates, that is the $\mathsf{XOR}$ of a linear function that takes a number of bits and the majority function for independent bits (or respectively a threshold function, that outputs $0$ if and only if the number of bits equal. to $1$ is less than a given threshold).
%\yr{Pas sûr que ce qui vient d'être dit va ici} \cb{Si, moi je trouve cela bien, mais il faut donner les formules pour ces fonctions, tout le monde ne les a pas en tête}.

%More formally, let $\mathsf{w}$ denote the Hamming weight, then the majority function on $t\in\mathbb{N}^*$ variables is defined as
%$$
%\begin{array}{ccccc}
%\mathsf{MAJ} & : & \{0,1\}^t & \to & \{0,1\} \\
% & & x & \mapsto & 0\text{ iff }\mathsf{w}(x) \leq \frac{t}{2}\,. \\
%\end{array}
%$$
%A threshold function on $t$ variables with parameter $\mathsf{th}\leq t+1$ is defined as 
%$$
%\begin{array}{ccccc}
%\mathsf{THR} & : & \{0,1\}^t & \to & \{0,1\} \\
% & & x & \mapsto & 0\text{ iff }\mathsf{w}(x) < \mathsf{th}\,. \\
%\end{array}
%$$

%\yr{En fait on peut définir que THreshold et dire que majorité c'est simplement le cas particulier $t/2$}


\subsection{Generic Cryptanalysis}
In~\cite{DBLP:journals/jar/AlekhnovichHI05,TCC:CEMT09,DBLP:conf/csr/ItsyksonS11,DBLP:journals/toct/CookEMT14,DBLP:journals/mst/Itsykson14} it was shown that Goldreich's one-way function was secure against some class of backtracking algorithms. 
An other generic cryptanalysis (that is without taking into considerations the specificities of the predicate) has been made by Bogdanov and Qiao~\cite{DBLP:conf/approx/BogdanovQ09} for linear stretch (that is $m=\lambda n$ for some constant $\lambda$), where they show that one cryptanalyst wins when it finds an $x'$ that is close with respect to the Hamming distance to the secret value $x$. The work then led to a subexponential attack~\cite{EPRINT:Applebaum15} in time $\mathsf{poly}(n)\cdot 2^{2\varepsilon n}$ if $m\in\Omega (n/ \varepsilon^{2d})$.

Naturally, the choice of the predicate should be non-degenerate: it should at least be unbiased and not linear. In the next section we describe some other criteria that the predicate should satisfy. 

More precisely, two classes of attacks arise when the predicate can be approximated by a linear function with fewer variables than $d(n)$ or when the degree or algebraic immunity of the predicate is small.

\subsection{Linearity Tests and Resiliency}
The use of $P_5$ predicate was shown to be secure against linearity tests first in~\cite{FOCS:MosShpTre03} up to a stretch $\mathsf{s} = 1.25-\varepsilon$ for any arbitrary $\varepsilon > 0$. Some other arguments were proven for non-degenerate predicates in~\cite{TCC:AppBogRos12}. Later, in~\cite{DBLP:conf/coco/ODonnellW14} it was shown to be secure in this class of attacks up to $\mathsf{s} = 1.5-\varepsilon$, also for arbitrary positive $\varepsilon$.
 
In the same paper~\cite{DBLP:conf/coco/ODonnellW14}, O'Donnell and Witmer also explained that if the predicate can be approximated with a function that takes $c<d$ variables as input, then the PRG can be broken as soon as $m\in\Omega(n^{c/2} + n\log n)$. This has been further analyzed in~\cite{EPRINT:Applebaum15}.

Naturally, this attack is thus naturally efficient when the predicate has a low resiliency order~\cite{DBLP:journals/tit/Siegenthaler84} which defines exactly the criteria: a Boolean function is $t$-resilient if it cannot be approximated by any boolean function with $t$ variables (the predicate $P_5$ has a resiliency order of $t=2$, and in the attack described above in~\cite{DBLP:conf/coco/ODonnellW14} uses thus an approximation of the predicate with $c=t+1$ variables, leading to an attack when the output size exceeds $n^{1.5}$.

\subsection{Algebraic Attacks}
Unsurprisingly, one of the most powerfull attack on Goldreich's PRG is the second type of attacks that is algebraic. Naturally, one can easily notice that a polynomial time algorithm inverts Goldreich's PRG when the output is in $\Omega(n^\delta)$ where $\delta$ is the algebraic degree of the predicate $P$. However, some other tricks can be done. 

A first example is in~\cite{STOC:AppLov16} where one important criteria to resist linear tests is in fact the \textit{bit-fixing degree}, that is the degree achieved by restricting the predicate $P$ by fixing some bits. An other criteria is also needed and explained in the same work: the algebraic immunity~\cite{C:Courtois03,EC:CouMei03} of the predicate that captures equations of smaller degree that can be derived from an equation of higher degree.

The output of Goldreich's PRG is \textit{a priori} in $\{0,1\}$ that is the field with two elements, but it can naturally be defined other a bigger alphabet $\mathbb{Z}_q$. In 2023, Akin Ünal proposed three distinguishing attacks~\cite{EC:Unal23} that work with different regimes. More precisely, if $m = n^{1+e}$, then there is an attack with space and time complexity in $n^{\mathcal{O}(n^{1-\frac{e}{\delta-1}})}$ with good advantage. A second attack is proposed that works when the alphabet is of size $n^{1-\frac{e}{\delta-1}}$. Eventually a third attack is proposed when the locality is $d$ and the field size $q$ is constant of complexity $2^{\mathcal{O}(n^{1-\frac{e'}{(q-1)d-1}})}$, where $e'<e$ and changing accordingly the advantage of the distinguisher.

Eventually and building up on its work, Akin proposed a way to extend the outputs of $\mathbb{Z}_2$ into larger fields in order to find algebraic relations in the outputs. Indeed, all previous attacks work only other larger fields. Up to now, its work has not been published and is only on eprint~\cite{EPRINT:Unal23b}.
\yr{Ce papier est à sur eprint depuis deux ans: un peu de mal à le traiter comme si c'était tout ok... à décider si on laisse ce paragraphe ou pas.}


\subsection{Instantiations and Concrete Cryptanalysis}
\yr{Cette section peut être divisée en deux ?}
Up to now, all the works that are described above give insights on whether a polynomial algorithm exist or not, and give criteria for the predicates for which the construction is good or bad. As we said before, many constructions are nowadays looking for specific applications and are looking to instantiate with concrete parameters such constructions. Therefore we need analysis that tackle this challenge, that will at the end provide confidence in instantiations for a very specific security level (e.g. 128 bits of security).

The first (but slightly different) construction that look like Goldreich's PRG is the FLIP family of stream ciphers~\cite{EC:MJSC16} and it's follow ups~\cite{DBLP:conf/indocrypt/MeauxCJS19} FiLIP~\cite{DBLP:conf/indocrypt/HoffmannMR20} and Elisabeth~\cite{DBLP:conf/indocrypt/HoffmannMS23}. However, those proposals suffer from some flaws~\cite{C:DuvLalRot16, AC:GBJR23} (mainly algebraic) that require to increase the size of the initial register to achieve a sufficient security level.

For a proposal without fixed parameters, it is very unlikely that a guess-and-determine approach will lead to a polynomial-time attack, it however reveals very powerful for subexponential-time attack and even more important it helps drastically to improve the complexity of attacks when the parameters are fixed. This has been done in~\cite{AC:CDMRR18} where the determine part is just a Gaussian elimination, together with some experiments that use Gröbner basis methods, but without the ability to prove or disprove the security with cryptographic sizes, as it is often the case in algebraic-like attacks. 

Four years later, the above cryptanalysis has been improved in~\cite{DBLP:journals/tit/YangGJL22} by making adaptative guesses, reducing the number of necessary guesses drastically and as the authors follow~\cite{AC:CDMRR18}, they broke the parameters sets intially proposed (for the $P_5$ predicate) and provide new ones as a challenge. Up to now, their parameters are not broken. In their work, Yang, Guo, Johansson and Lentmaier also propose a very interesting strategy: the guess-and-decode. Eventually, their work also targets other predicates but without an analysis of the algorithm, as each predicate requires an other deep work of analysis.

\subsection{Bit-Fixing Correlation Attack}
While resiliency and algebraic immunity appeared to be the two most important criteria, it appeared recently that those ones are not enough~\cite{EPRINT:FLLL24}. Indeed, for \textsf{XOR-MAJ} and \textsf{XOR-THR} functions, their resiliency order and algebraic immunity is good, but in~\cite{EPRINT:FLLL24}, the authors observed that if some bits are guessed to be zero's or one's, then the output can be seen as noisy linear equations, with a noise much higher than the original one that is given by the original function. The attack is powerfull as the number of bits taken in entry by the majority or threshold function increases. Eventually this work highlights again the necessity for having dedicated cryptanalysis and asks the natural question of a good predicate that will be secure against algebraic attacks, linearity tests and naturally guess-and-determine strategies.

In Table~\ref{tab:subexpattacks} we give some subexponential attacks where $d$ denotes the locality of the PRG, $\mathsf{s}$ is the stretch, that is the output size is in $n^{\mathsf{s}}$ and $n$ is naturally the size of the seed.

\begin{table}[!ht]
\centering
\begin{tabular}{lcc}
\toprule
Technique & Complexity & Reference\\
\midrule
Test and find $x'$ close to $x$ & $2^{\mathcal{O}(n^{1-(\mathsf{s}-1)/2d})}$ & \cite{EPRINT:Applebaum15,DBLP:conf/approx/BogdanovQ09} \\
Guess-and-determine & $2^{\mathcal{O}(n^{2-\mathsf{s}})}$ & \cite{AC:CDMRR18} \\
Adaptative Guess-and-determine & $\mathcal{O}(n^2 2^{\frac{n^{2-\mathsf{s}}}{4}})$ & \cite{DBLP:journals/tit/YangGJL22} \\
Guess-and-decode & $\mathcal{O}(n^\mathsf{s} 2^{\frac{n^{2-\mathsf{s}}}{4}}$ & \cite{DBLP:journals/tit/YangGJL22} \\
\bottomrule
\end{tabular}
\caption{\label{tab:subexpattacks}Subexponential and concrete cryptanalysis. The first technique applies for any predicate of locality $d$ while the three other look into the $P_5$ predicate.}
\end{table}







\subsection{Targets for cryptanalysis}
The high versatility in Goldreich's PRG asks many questions. First it is not clear how to choose the best predicate that will achieve the best stretch $\mathsf{s}$ for the smallest locality $d(n)$. 

Maybe the simplest target would be the instantiation with the $P_5$ predicate and which are given in~\cite{DBLP:journals/tit/YangGJL22}.
\begin{table}[!ht]
\centering
\begin{tabular}{lc}
\toprule
Seed size & stretch \\
\midrule
1024 & 1.02  \\
2048 & 1.10 \\
4096 & 1.19 \\
\bottomrule
\end{tabular}
\label{tab:paramgar}
\caption{Challenge parameters for Goldreich's PRG with $P_5$ predicate for a security level of 128 bits.}
\end{table}

\subsection{Implementation}
The PRG differs from the other weak shallow PRFs presented, in the sense that we do not have a key. However, the adversary is given the subset $\sigma_i$ that are uniformly distributed among the sets of subsets of size $d$ and this plays the role as the $x$'s in the other algorithm (it is public but not choosable by the attacker). The secret is the seed that is called previously $x$ as the seed and for readability we call it $k$ in Algorithm~\ref{alg:gar}.

\yr{Je veux bien qu'on appelle $k,x$ comme les autres j'essaie un truc mais je suis pas si convaincu. J'utiliserais au moins $\sigma$ ?}

\begin{algorithm}
  \caption{\label{alg:gar}The Goldreich's PRF/PRG \\
    Parameters: $n$ (seed size), $m$ (output size), $d$ (locality) and $P:\{0,1\}^d \to \{0,1\}$ a predicate.}
  \begin{algorithmic}
  \State \textbf{Input:} $k\in\{0,1\}^n$, $x\in [1,\ldots,A^n_d]^{m}$
    \State $y \gets \epsilon$
    \For{$i = 1$ to $m$}
    	\State $y \gets y||P(k[x_i])$
	\EndFor
    \Return $y$
  \end{algorithmic}
\end{algorithm}




%%% Local Variables:
%%% mode: latex
%%% ispell-local-dictionary: "english"
%%% TeX-master: "main"
%%% End:

\input{vdlpn}
\subsection{EALPN}
\label{sec:ealpn}


\subsubsection*{Original Reference}
Boyle, Couteau, Gilboa, Ishai, Kohl, Resch, Scholl 
\cite{C:BCGIKR22} \cite{C:RagRinTan23}

\subsubsection*{Security Arguments and Cryptanalysis}
%!TODO! Write about security arguments and cryptanalysis


\subsubsection*{Parameters}
% !TODO! write about the parameters and their role


\subsubsection*{Implementation}
%!TODO! write about the implementation 

Algorithm~\ref{alg:ealpn}

\begin{algorithm}
  \caption{\label{alg:ealpn}The {\tt EA-LPN} PRF \\
    Parameters:  integers $N$, $t$, $\ell$,  $w = \lceil \frac{5N}{t} \rceil$}
  \begin{algorithmic}
  \State \textbf{Input:}  $x = ((j_1,x_1), \dots, (j_\ell,x_\ell)) \in ([0, \frac{t}{\ell} - 1] \times \mathbb{Z}_{2^w})^\ell$; $k = (k_0, \dots, k_{t-1}) \in (\mathbb{Z}_{2^w})^t$
    \State $r \gets 0$
\For{$i = 1$ to $\ell$}
    \If{$x_i \geq k_{j_i}$}
        \State $r \gets r + 1$
    \EndIf
\EndFor\\
\Return $r \bmod 2$
 \end{algorithmic}
\end{algorithm}



%%% Local Variables:
%%% mode: latex
%%% ispell-local-dictionary: "english"
%%% TeX-master: "main"
%%% End:

\section{Conclusion}
\label{sec:conclusion}

\begin{table}
  \centering
  {
    \renewcommand\arraystretch{1.4}
    \setlength\tabcolsep{0.4em}
    \footnotesize
    \begin{tabular}{ccc|ccc|cc}
      \toprule
      \multicolumn{3}{c|}{Description} & \multicolumn{3}{c|}{Spaces}  & \multirow{2}{*}{NL op.} &  \multirow{2}{*}{Best attack} \\
      Name & Use Case & Ref. & Input & Key & Output \\
      \midrule
      % Alternating Moduli
      \multirow{2}{*}{Alternating} & \multirow{3}{*}{MPC}
      & \cite{TCC:BIPSW18} & $\Z_2^{384}$\dag & $\Z_2^{384}$\dag & $\Z_2$ & $\mod$ & ?? \\
      \cline{3-8}
      \multirow{2}{*}{Moduli} & &  \multirow{2}{*}{\cite{C:DGHIKS21}} & $\Z_2^{256}$\ddag & $(\Z_2^{256})^{256}$\ddag & \multirow{2}{*}{$\Z_2^{81}$} & \multirow{2}{*}{$\mod$} & ?? \\
      & &  & $\Z_2^{320}$ & $(\Z_2^{320})^{320}$ & & & ?? \\
      \midrule
      % Goldreich
      \multirow{1}{*}{Goldreich's} & \multirow{2}{*}{??} & \multirow{3}{*}{\cite{DBLP:journals/tit/YangGJL22}} & $(A_5^{1024})^{1176}$ & $\{0,1\}^{1024}$ & $\{0,1\}^{1176}$ & \multirow{3}{*}{$P_5$} & ?? \\ % m=n^1.02
      \multirow{1}{*}{PRG} & & 
      & $(A_5^{2048})^{4389}$ & $\{0,1\}^{2048}$ & $\{0,1\}^{4389}$ & & ?? \\ % m=n^1.10
      & & & $(A_5^{4096})^{19893}$ & $\{0,1\}^{4096}$ & $\{0,1\}^{19893}$ & & ?? \\ % m=n^1.19
      \cline{3-8}
      & & \cite{EC:ABBS25} & $(A_5^{8192})^{45387}$ & $\{0,1\}^{8192}$ & $\{0,1\}^{45387}$ & $P_5$ & ?? \\ % m=n^1.19
      \bottomrule
    \end{tabular}
  }
  \caption{A summary of the wPRFs we present. The ``\dag'' symbol means the primitive is now considered unsafe, while ``\ddag'' indicates an ``aggressive'' set of parameters. The data limit for all such functions is at least $2^{40}$, and the intended security level at least 128 bits. ``NL op.'' refers to the non-linear operation.}
  \label{tab:summary}
\end{table}


The goal of this Systematization of Knowledge (SoK) article was to unify and present four distinct families of weak pseudo-random functions (wPRFs) under a common lens. These families of wPRFs have emerged from the theoretical cryptography community in response to various application needs. Despite their differing origins, all the constructions we examined share a common characteristic: they are shallow, typically consisting of a single round of computation. Moreover, they operate within weakened security models when compared to standard PRFs or PRPs, notably by restricting the attacker from choosing inputs. More importantly, the inputs to these primitives  are sampled uniformly at random.

This uniform randomness of public inputs, combined with the small domain sizes typically proposed, makes differential cryptanalysis largely impractical in this context. As a result, the main attack strategies that remain applicable seem to be linear cryptanalysis, based on exploiting statistical biases in (linear combinations of) the output, and algebraic or combinatorial techniques. However, the field currently lacks a comprehensive understanding of how effective such attacks can be in practice.

In particular, a more comprehensive combinatorial analysis is needed to characterize the conditions under which an attack may succeed based on properties of the public inputs. Indeed, attacks may become feasible when public inputs fall within certain ``bad'' subsets. Identifying and understanding the structure of these subsets is an interesting and important research direction.

More broadly, there is a remarkable lack of systematic cryptanalysis in this weak-function setting. Security claims for these constructions often rely on heuristic or partial evidence, and the distinction between exponential and polynomial-time attacks becomes blurred once concrete parameters are fixed. Despite this, many of these primitives are being integrated into advanced cryptographic protocols and proposed for real-world deployment. We therefore encourage the symmetric cryptography community to analyze these constructions, and aim to facilitate this task by providing reference implementations for each of them. Targeted and rigorous cryptanalysis is essential before these primitives can be confidently used inside high-level protocols. Bridging the gap between theoretical proposals and concrete security evaluation is not just desirable,  it is necessary to ensure the security of the advanced protocols that build upon them.

%\begin{itemize}
%\item On a besoin de quantifier précisément la complexité des algorithmes d'attaque histoire de pouvoir dimensionner correctement et précisément les instances sensées offrir un niveau de sécurité donné
%\item en l'état, il n'y a pas de cryptanalyse dans le modèle ``weak nonce'' ou ``weak IV'' (demander à Anne ce qui existe), qu'on introduit
%\end{itemize}


%All constructions that we presented consist of a weakened security model than the standard one (PRF or PRP), as an attacker is not able to choose the public inputs. It is however somehow different than a known plaintext scenario, as the inputs here are taken uniformly at random.

%The fact that public inputs are taken uniformly at random, combined with the proposed size makes the differential cryptanalysis unpractical. What remains seem to be, up to now likewise linear cryptanalysis, by observing biases in (linear combination of) the output or algebraic-like attacks. \yr{C'est les seules types d'attaques partout non ?}

%Some combinatorial analysis is also missing to identify a success probability of an attack that is computable over the public inputs. Indeed, as Kerckhoff's principle states that the security should rely only on the secret of the key, this principle is not met in those constructions as it could appear than an attack could work when the public inputs lie in a specific subsets. The identification of ``bad'' public input set is thus a very interesting question that need to be answered.
%\yr{Faut-il réinsister plus ?}

%Eventually, we believe that this weak scenario drastically lacks of security analysis. Indeed, actual security relies on cryptanalysis. We believe this to not be fulfilled, specifically because when fixing parameters, the distinction between exponential/polynomial types of attacks do not make anymore sense. Nowadays, many of those proposals have a purpose to be used in practice in many advanced cryptographic protocols. Therefore a big gap is present and we think that the cryptographic community should collectively target those constructions, with careful analysis in order to gain confidence in advanced protocols before proposing them as being used and implemented.


\bibliographystyle{alpha}
\bibliography{biblio,abbrev3,crypto}

\end{document}
